% 本文件由 LaTeX 排版 Agent 根据有效 translation.json 自动生成。
% author=Gregory of Nyssa; work_id=2912; title=论婴儿夭折

\chapter{论婴儿夭折}

% structure: p0001 source=h1 target=omitted reason=page-title-already-rendered-by-parent

每一个论文作家和小册子作者，至为卓越的阁下，都会以你为题材来展示其辞藻；你的奇妙品质是一处广阔的跑马场，供他驰骋。一个高尚且富有启发性的题材，在能手中确实能造就更雄伟的风格，将其提升至伟大现实的高度。然而，我们如同一匹老马，将留在这一提议的跑马场之外，只侧耳倾听那称颂你赞词的竞赛，若有任何文学战车的喧嚣全力驶过这些奇观，传到我们耳中。但是，尽管年老可能迫使一匹马远离赛场，却常常发生这样的情况：奔腾赛马的喧嚣会激起它的兴奋，它热切地抬起头，在呼吸中流露精神，频频跃起并刨地，尽管这种热切已是它所剩的一切，而时间已耗尽了它奔跑的力量。同样，我们的笔也留在竞赛之外，年老迫使它将赛场让给当下兴盛的学者；然而，它对你参与竞赛的热望依然存活，它仍能展现这一点，即使当下兴盛的那些文体家正处于能力的巅峰。但所有这些我对你热忱的展现，都与赞美你本人无关：没有哪种风格，无论多么有力且均衡，能轻易在其中成功；因此，任何试图描述你品格中那种令人尴尬却又和谐的对立混合的人，都必然远远落后于你的真实价值。的确，自然通过在前睫毛的阴影下射出耀眼光芒之前，给眼睛本身带来较弱的光线，使我们能够忍受日光，因为它按我们所需的比例与睫毛投下的阴影混合。同样，你品格的宏伟与伟大，被你的谦逊和谦卑之心所调和，不仅不刺眼，反而使视线愉悦；其中，谦卑之心并不使伟大的光辉暗淡，也不使其潜在的力量被忽视；反而，在谦卑中可见其高贵，在高贵中可见其谦逊。这些须由他人来描述；此外，还要颂扬你心智的多重洞见。你的智慧之眼，或许可以说，多如头发；它们敏锐无误的目光同样注视着一切；远见未来；近处不疏忽；它们不等待经验来教导权宜；它们用希望的洞察，或用记忆的洞察来看；它们全面审视当下；先在一事，再在另一事，但不混淆，你的心智以同样的精力并按所需的专注工作。另一个人也必须记录他对你使贫穷变为富足的方式的赞叹；如果在我们这个时代确实能找到有人将此作为赞美和惊奇的对象。然而，现在，如果以前从未有过，对贫穷的爱将通过你而丰盛，你的\Emph{\OriginalTerm{未得之财}{ingotten}}将比\Person{克罗伊斯}的\Emph{金条}更受羡慕。因为海上和陆地，连同其自然产物的全部嫁妆，曾使谁富足，如同你拒绝世间的富足使你富足那样？他们清除钢上的污渍使其如银般光亮：同样，你生命的光辉变得更加明亮，因为被财富的锈蚀精心擦净。我们将这些留给那些能详述的人，也留给他们详述你对那些获得比不获得更荣耀之事的卓越认识。但请允许我说，你并非藐视所有获得；有些获得，虽然你的任何前任都未能抓住，但唯独你双手抓住了；因为，代替衣袍、奴隶和金钱，你拥有并持守人的灵魂，并将它们储存在你爱的宝库里。那些因这类赞美而获得荣耀的论文作家和小册子作者，会探讨这些事情。但我们的笔，如今已是老手，只应振作到足以以步行速度通过你所提出的智慧问题；即：我们应该如何看待那些过早被带走的人，其出生时刻几乎与死亡时刻重合？有教养的异教\Person{柏拉图}曾假借一个复活之人的口，谈论了许多关于另一世界审判法庭的哲学；但他将这个问题留作一个奥秘，似乎过于伟大以至于不能用人间的推测来探讨。那么，如果我们这些辛勤思考中有任何能够澄清此问题晦涩之处的，你无疑会欢迎这一新的解释；否则，你至少会原谅这老迈之人的尝试，并接受，如果没有别的，我们想给你带来些许愉悦的愿望。历史记载，\Person{薛西斯}，那位伟大的君王，几乎使天下万国成为一个大营，并以他自己的计划唤醒了整个世界，当他征讨希腊人时，愉快地接受了一个穷人的礼物；那礼物是水，并且不是装在罐中，而是捧在手心里。所以，你以天生的慷慨效法他；对他而言，心意使礼物有价值，而我们的礼物本身可能只是贫乏的水。对于天上的奇观，一个人看见它们的美丽，无论他是受过训练观察它们，还是以非科学的眼光仰望，都是同等的；但对它们的感觉，对于一个从哲学出发来沉思的人，和一个只凭感官知觉来接受的人，并不相同；后者可能喜悦阳光，或认为星辰的美丽值得赞叹，或观察了月亮在一个月中的阶段；但前者，拥有\OriginalTerm{灵魂洞察}{soul\-insight}，并且训练有素使他能够理解天象，他忽略了那些取悦较无思想者感官的事物，而着眼于整体的和谐，审视那些由相反运动在圆形旋转中产生的协和；内圈如何以与恒星旋转相反的方向转动；在内圈中观察到的天体如何在其接近和远离、相互遮挡和侧翼运动中形成各种组合，却始终以完全相同的方式产生那著名且永恒的和谐；那些不忽视最微小星星位置、且其心智因训练而居于上界的人，能意识到这一点，他们同样关注所有星星。同样，你，我宝贵的生命，注视着\OriginalTerm{天主经纶}{Divine economy}；你将那些不断占据群众心智的物体——我指的是财富、奢侈和虚荣——如阳光在脸上闪烁，使无思想者目眩——置于一旁，而不忽略世上看似最微不足道的问题；因为你确实最仔细地审视人类生命中的不平等；不仅关于财富与贫困，以及地位和出身的差异（因为你明白它们都是虚无，其存在不是由于任何内在真实，而是由于那些被虚无如实物般击中的人的愚蠢估计；并且，如果只从某个因荣耀而闪耀的人那里抽走盲目崇拜者的仰慕，那么在他那使他得意的膨胀骄傲之后，即使全世界的财富都埋在他的地窖里，也将一无所有），但你的焦虑之一是知道，在神圣管理的每个细节的其他意图之中，为何有的人生命延长到老年，而另一个人只得到如此一部分，以至于仅吸一口气就死去。如果世上没有一件事不是出于天主，而是所有都连接于天主的旨意，且如果神性是有技巧和审慎的，那么必然在这些事中有某种计划，带着祂的智慧，同时也带着祂的眷顾的印记。盲目的无意义事件绝不可能是天主的工作；因为，正如圣经所说，\BibleQuote{以智慧造成万物}，是天主的属性。那么，在以下情况下我们能辨别出什么智慧呢？一个人进入生命的舞台，吸进空气，以痛苦的哭喊开始生存的过程，向自然付出了眼泪的贡品，在尝到任何甘甜之前，在感觉有任何力量之前，就尝到了生命的痛苦；关节仍然松弛，柔软，未成形；一句话，在他甚至成为人之前（如果理性的恩赐是人的特征，而他从未拥有它），这样的人，在母腹中的胚胎相比只有见过空气的优势，如此短命，死去并又分解；或者被抛弃，或窒息，或出于自身的软弱而停止生存。我们对此该作何想？我们对这样的死该如何感受？这样的灵魂会见到它的审判者吗？它会与其他灵魂一同站在审判台前吗？它会因一生所行之事接受审判吗？它会根据福音的言论，被火净化，或被祝福的露水清凉，而得到公正的报应吗？但我看不出我们如何能想象，对于一个这样的灵魂，会如此。报应一词暗示之前必须给出某种东西；但根本没有生活过的人已被剥夺了给出任何东西的材料。因此，既然没有报应，就没有善恶可期待。报应意谓着要回报这两种性质之一；但既不在善的范畴也不在恶的范畴中的事物，根本不在任何范畴中；因为善恶之间的对立是排他的，不容中间；对没有开始任何一方的人，两者都不会来。因此，落在两者之下的东西，可以说甚至不曾存在。但如果有人说，这样的生命不仅存在，而且作为善的生命之一存在，并且天主赐予（而非偿还）善给这样的人，我们可以问，他提出这种偏爱的理由是什么；在这样的观点中，公义如何显明；他如何证明他的想法与福音的言论一致？在那里，主说，\BibleQuote{你们做了这些事，然后得王国为赏报是应当的。}但在此例中，没有先行的行动或意愿，因此有什么理由说这些人会从天主那里得到任何预期的报偿？如果有人毫无保留地接受这样的说法，即任何这样进入生命的人必然被归类于善者，那么他就会意识到，根本不参与生命比活着更幸福，因为在一例中，即使出身于蛮族父母或不合法的结合，善的享受也毫无疑问；但活了自然可能达到的年限的人，其生命必然或多或少混有罪恶的污染，或者，如果要完全脱离这种感染，则要付出许多痛苦的代价。因为德行并非没有斗争就能为其寻求者所获得；对人性而言，戒绝享乐之路并非无痛的过程。因此，有两种考验是寿命较长者的必然命运：要么在此世德行的艰辛场地上战斗，要么在来世为邪恶的一生承受痛苦的报应。但在早夭的婴儿身上，没有这类事情；他们立即通过到有福的命运，如果持此观点的人所言不虚。由此也必然推出，无理智状态比有理智更可取，德行因而显示为毫无价值：如果从未拥有德行的人不受损失，那么，就享受福乐而言，为获得德行而劳动便是无用的愚蠢；无思想的状态将成为从天主审判中得胜的一方。出于这些及类似理由，你吩咐我仔细探究此事，以期通过严谨推理的探究，得到某种坚实根基，使我们的思想得以安立。\par

\BibleQuote{啊，天主的智慧和知识的财富多么深奥！他的判断多么不可探测，他的道路多么不可追寻！因为谁认识主的心呢？\BibleRef{罗 11:33}}\newline{}但另一方面，同一宗徒宣称，属于神性的人的特质是\BibleQuote{判断一切事，\BibleRef{格前 2:15}}，并称赞那些\Quote{得著了丰富\BibleRef{格前 1:5}，藉着天主的恩宠，\InnerQuote{在一切言语和一切知识上，}}我敢断言，不可忽略我们能力范围内的考察，也不可对所提出的问题不作任何探究或没有想法；否则，就像我们当前讨论的主题一样，这篇论文可能会以无果而终，在成熟之前就被那些不愿鼓起勇气寻求真理的人的致命怠惰所破坏，如同新生儿尚未见到光明、尚未获得力量一样。我也断言，不要像在法庭上辩护那样立即反驳和应付反对意见，而是要在讨论中引入一定的秩序，并一步步引导观点。那么，这个秩序应当是什么？首先，我们想知道人性的来源，以及它为何会存在。如果我们找到了这些问题的答案，就不会得不到所需的解释。现在，在智力或感官世界中，除了天主以外的一切存在物都归因于祂，这是一个无需证明的命题；任何对事物真相稍有了解的人都不会否认它。因为每个人都同意宇宙与一个第一因相连；其中没有任何东西是自存而成为自己的起源和原因的；相反，只有一个非受造的永恒本质，永恒不变，超越我们对距离的一切概念，被认为无增无减，超出任何定义的范围；时间、空间及其一切结果，以及在这些之前思想所能把握的、可理解的超世俗世界中的任何事物，都是这个本质的产物。那么，我们断言人性是这些产物之一；圣经的教导也帮助我们，它宣告：当天主将其他一切事物带到生命舞台之后，人出现在地上，是神圣来源的混合物，神明般的理智本性与他形成所需的几部分尘世元素结合，被造物主塑造为神圣超越能力的化身肖像。不过，最好直接引用原话：\BibleQuote{天主于是照自己的肖像造了人，就是照天主的肖像造了人。\BibleRef{创 1:27}}现在，这个活物的创造原因，我们以前的某些作者已经作如下说明。整个受造物分为两部分：宗徒所说的\Quote{所见的}和\Quote{所不见的}（后者指可理解的、非物质的，前者指可感觉的、物质的）；这样分开之后，属于\Quote{所不见的}天使和神性本性，居于世界之上、诸天之上，因为这样的居所与他们的构造相符；因为\OriginalTerm{理智本性}{intellectual nature}是一种精微、清晰、无碍、敏捷的东西，而天体是精微、轻盈、永远运动的；相反，居于可感觉事物最末端的尘世，绝不可能是理智受造物适宜居住的地方。轻飘的与沉重的之间能有什么对应呢？为了使大地不完全缺乏理智和非物质的本地位居，人（这些作者告诉我们）是由至高上智所塑造，其尘世部分包裹在灵魂的理智和神明本质之上；因此，这种与物质重量的结合使得灵魂能够生活在具有某种与血肉之实质亲缘关系的尘世元素上。那么，所有正在诞生的东西的设计，就是使高于诸天和尘世宇宙的大能，能够通过理智本性的部分在受造物的各部分受到光荣，这些本性因同一运行能力（即注视天主的能力）而共同趋向同一目的。但注视天主的这一运行，恰如其分地说，是理智本性同类的\OriginalTerm{生命滋养}{life-nourishment}。正如这些尘世的肉体由尘世的滋养所保存，我们在它们身上都看到物质性活力的运行，无论是动物还是理性生物，同样有理由假定也存在一种理智的\OriginalTerm{生命滋养}{life-nourishment}，使这类本性得以维持存在。但如果身体的食物，虽然来来去去、循环不已，仍能给获得它的人带来一定的生命力，那么，\OriginalTerm{参与天主}{partaking of God}那永远不变、始终如一的事物，岂不是更会使食用者得以保存？那么，如果这就是理智本性的\OriginalTerm{生命滋养}{life-nourishment}，即\OriginalTerm{参与天主}{partaking of God}，这分享不能由相反性质的东西获得；想要分享的人必须在某种程度上与所分享之物相似。眼睛因本身有光来捕捉同类之光而享受光明，手指或任何其他肢体不能产生视觉，因为其中没有这种自然光以有组织的方式存在。同样的必要性要求，在我们\OriginalTerm{参与天主}{partaking of God}的行动中，分享者的构造应与所分享之物有某种亲缘关系。因此，正如圣经所说，人是按\OriginalTerm{天主的肖像}{image of God}被造的；我认为，这使同类能看见同类；而看见天主，正如前面所说，就是灵魂的生命。但既然对真善的无知如同迷雾遮蔽灵魂的视觉敏锐力，当这迷雾变得越来越浓，形成如此厚重的云层，以致真理之光无法穿透这无知的深渊，那么，随着光明的完全剥夺，灵魂的生命完全停止；因为我们说过，灵魂的真正生命在于参与善；但当无知阻碍对天主的认识，灵魂因而停止\OriginalTerm{参与天主}{partaking of God}，也就停止活着。但没有人能强迫我们给出这无知的家族史，问它从何而来、父亲是谁；让他从字本身理解，\Quote{无知}和\Quote{知识}表明灵魂的一种关系；但任何关系，无论是否表达，都不传达实体的概念；关系和实体是截然不同的类别。那么，如果知识不是实体，而是灵魂的一种完善运行，就必须承认无知离实体更遥不可及；而不属于实体范畴的东西根本就不存在；因此，我们不必费心追问它从何而来。然而，圣言宣告，活于天主就是灵魂的生命，而这一活于天主就是按照各人的能力认识天主，无知并不表示任何实体的真实性，而只是认识运行的否定，并且一旦不再实现与天主的参与，灵魂的生命立即被取消，这是最坏的事——因此，一切善的创造者会在我们身上运作对这种恶的治疗。治疗是好事，但一个不关注福音奥迹的人仍会不知道治疗的方式。我们已经表明，与生命之主天主分离是一种恶；那么，这种病弱的治疗就是重新与天主为友，从而重新拥有生命。既然这样的生命总是在希望中向人类展示，就不能说获得这生命绝对是对好生活的奖赏，而其反面是对坏生活的惩罚；但我们所坚持的类似于眼睛的情况。我们不说，有清晰视力的人因能看到视觉对象而得到奖赏；也不说，眼睛有病的人因某种刑罚而经历视觉功能的丧失。眼睛在自然状态下，视觉必然随之而来；被疾病损坏时，视力的丧失也必然随之而来。同样，真福的生命对于那些保持灵魂官能清晰的人来说，如同熟悉的第二天性；但当无知盲目之流阻止我们参与真光时，我们就必然失去那享受，而享受本身我们称之为参与者的生命。\par

如今我们既已奠定了这些前提，是时候凭此来审视向我们提出的问题了。问题大致如下：\Quote{如果真福的赏报是按正义的原则分施的，那么一个在襁褓中夭折、此生既未奠定任何或善或恶的基础、以致无法按其功过获得任何回报的人，应被置于哪一类中？}对此，我们将着眼于我们已经阐明的结论来作答：这未来的福乐，在其本质上虽是人类的天赋产业，却也可以在某意义上被称为赏报；我们将用与先前相同的实例来阐明我们的意思。假设有两人患了眼疾：一人极其勤勉地接受治疗，忍受医学所能施予的一切，即使痛苦异常；另一人则毫无节制地沉溺于沐浴和饮酒，对医生关于治眼的任何忠告都充耳不闻。那么，当我们看到这两人的结局时，我们说他们都各得其选择之果：一人失去光明，另一人享有光明；我们误用词汇，竟将这必然的结果称为赏报。同样，关于婴儿的问题，我们也可以这样说：未来生命的享受，确实属于人类的权利，但鉴于无知之灾已几乎攫取了所有现今生活在肉身中的人，那藉必要的治疗程序洗净了自己的人，在进入真正符合本性的生命时，便获得了他勤勉的应得赏报；而那拒绝德行之清泻剂、并因所陷溺的享乐而使其无知之灾难以治愈的人，便陷入了反乎本性的状态，从而与真正符合本性的生命疏离，无分于那本属于我们、与我们相宜的存在。至于无辜的婴儿，其灵魂之眼前并无这灾障遮蔽其光明之尺度，因此他继续存在于那本性的生命中；他不需要来自净化的健全，因为他从未将灾障接纳于灵魂之中。此外，在我看来，今世生命与我们所望的未来生命似乎提供了某种类比，且紧密相连，如下：最柔弱的婴儿以母乳喂养；然后，另一种适合这养育对象并密切适应其需要的食物接续而来，直到最后长成。同样，我认为，灵魂以不断适应其能力的渐进过程，以某种有规律的前进，分有那真正本性的生命；它按自己的容量和能力，接受真福境界的喜乐之量。事实上，我们从保禄那里也学到这一点，他对已在德行中成长的人和对不完全的\Quote{婴孩。}他对后者说：\Quote{我原是给你们喝奶，并不是给饭吃，因为那时你们还不能吃\BibleRef{格后 3:2}。}但对那些已达理智成熟之圆满尺度的人，他说：\Quote{唯独成年人的食粮，才是那些因习用而感官习练了的人…\BibleRef{希 5:14}}。然而，不能说那成人和那婴孩处于相同状态，尽管两者都免于任何疾病的接触（因为不分享完全相同的事物的人，怎能处于同等的享受状态呢？）；相反，尽管两者都免于疾病之苦可同等地断言，只要两者都不受其影响，但论到喜乐时，尽管感受者处于相同状态，享受却无相似之处。对于成人，有自然的喜乐在于讨论、处理事务、光荣地履行职务、因扶助贫者而显赫；或许在于与所爱的妻子共度生活、管理家务；以及今生一切消遣中的娱乐，如音乐作品、戏剧表演、狩猎、沐浴、体操、宴会的欢乐，以及其他诸如此类。而对于婴孩，则自然地以其奶水、保姆的怀抱、以及轻柔的摇晃使其入睡并使之香甜为乐。超越于此的任何福乐，其年幼的年岁自然地阻止他感受到。同样，那些在此生中藉德行之道滋养了灵魂力量的人，用宗徒的话来说，其心智的\Quote{感官}已\Quote{习练了，}若他们被转移到那超越的、脱出肉身的来世生命中，将按他们所达到的状态和能力，按比例分有那神圣的喜乐；他们将按其获得的能力，或多或少地享有其财富。但那从未尝过德行之味的灵魂，虽然确实可以完全免于来自邪恶的痛苦，因为从未染上恶疾，然而，在起初，它只是如这乳儿所能接受的那样，分有那来世生命（按照我们之前的定义，在于认识并存在于天主内）；直到它因瞻仰那真正存在者如同适合的滋养而茁壮成长，并变得能接受更多时，才随心从这真正存在者的丰盛供应中汲取更多。\par

因此，基于所有这些考虑，我们认为，一位在旅程中达致一切德行者的灵魂，与一位其生命之份完全空白者的灵魂，同样不受来自邪恶的苦难的影响。尽管如此，我们并不认为他们生活的运用处于同一等级。前者听到了那些属天的宣告，即，按先知的话：\Quote{天主的荣耀被宣告}，并游历受造界，被引导认识造化主宰；他以真智慧为师，即宇宙景象所提示的那智慧；当他观察物质的日光之美时，他借类比掌握了真阳光之美；他在这坚实的大地中看到其造物主的不变性；当他感知苍穹的广阔时，他走上通往那环绕宇宙之大能的无限广阔的道路；当他看到太阳的光芒从如此高远之处乃至到达我们时，他开始相信，通过此类现象，神性理智的活动不会不从天主性的高处降下，甚至到我们每个人；因为如果一个光体能以光的力量占据其下的一切事物，并且，除此之外，能在向所有能使用它的人提供自身的同时仍保持自足不分散，那么那光体的造物主岂不更成为\Quote{万物之中的万有}，如宗徒所说，并以祂自身那样的度量进入每个能接受祂影响的事物？不，且看麦穗、植物萌芽、成熟的葡萄、初秋之美（无论是果实还是花朵）、自发萌生的草、高耸入云的山峰、山坡上泉水从丰盈的胸怀涌出并流经峡谷成为河流、大海从四面八方接收河流却仍保持在界限内，波涛被沙滩边缘限制，从不跨越大陆的固定界限：看到这些和类似景象，理性的眼睛怎能不在其中找到我们对真实之物的教育所需的一切？一个欣赏这些景象的人，是否只为自己获取了享受那些超越之乐的微弱能力？更不用提那些磨砺心智趋向道德卓越的学问：几何学、天文学、数理科学所提供的真理知识，以及提供未知证明和已知确信的每种方法，而超越这一切的，是受默感著作中所含的哲学，它完全净化那些在其中受教育者，使之通达天主的奥秘。然而，对这些事物一无所知、甚至未由物质宇宙引导感知其上之美、带着一种柔嫩、未成形、未受训的心智度过一生的人，他不像我们的论证所表明的那另一个人那样会被置于同样的环境中；因此，据此观点，不能再坚持说，在两种假设的完全相反的情形下，未参与生命的人比高尚参与生命的人更有福。当然，与一生生活在罪中的人相比，不仅是无辜的婴孩，甚至从未来到世上的人也是有福的。我们从福音中关于犹达斯所宣判的句子也学到这一点（\BibleRef{玛 26:24}）；即，当我们想到这样的人时，不存在的东西比那样存在于此罪中的东西更可取。因为对于后者，由于根深蒂固之邪恶的深度，以炼净方式施行的惩罚将延至无穷；但对于从未存在的东西，怎样折磨能触碰它呢？——然而，尽管如此，将婴儿般不成熟的生命与完美德行进行比较的人，他本身应被判定为不成熟，竟以此判断真实。那么，你因此问，为什么某个年龄很小的人被悄然带走离开活人之地？你问上智在此默观什么？好吧，如果你想到所有那些作为非法结合之证据而被父母除掉的婴儿，你无权以此不义行为要求天主交代，祂必定审判如此行的不圣洁行为。另一方面，至于某个婴儿，虽父母养育了他，并以养育和祈祷热切照顾他，但他仍不在世上逗留，而是因疾病甚至死亡（这是其确切原因），我们冒昧提出以下考虑。这是天主眷顾完美的一个迹象：祂不仅医治已存在的病患，也确保一些东西从不参与祂所禁止的事；也就是说，有理由期待，预知未来如同知道过去的那位，应阻止一个婴儿发育至完全成熟，以便祂预知在其未来生命中已察觉的恶不致发展出来，并为了赐给一个终生有恶习的人的生命不会成为他恶习的实际材料。我们将用一个比喻更好地解释我们的想法。\par

假设有一场极其丰盛的宴会，为若干宾客预备，由他们中的一位担任主席，此人天赋能准确了解每一位的体质特点，以及何种食物最适合各人的性情，何种有害不宜；此外，他被赋予对他们近乎绝对的权威，可随意允许某人留在席间或驱逐某人，并采取一切预防措施，使每个人都食用最适合其体质的菜肴，好让病人不致因饮食助长其病而自取灭亡，而体质健壮的宾客也不致因食用了不适合他的东西而生病，或因过饱而感到不适。假设其中有一个好饮酒的人，在宴会中途甚至一开始就被带出去；或让他留到最后，全凭主席如何确保整个席间尽可能维持完美秩序，并避免暴食、酗酒和醉态的丑恶景象。正是如此，当那个人不甘心被从所有美味佳肴中拉走，失去他喜爱的美酒时，他往往会指责主席缺乏公正和判断，认为是出于嫉妒而非对他有任何预虑而将他逐出盛宴；但若他看到那些因长时间饮酒已经开始失态的人，呕吐、伏桌、言语不雅，他或许会感激主席在他陷入那种状态之前将他从深沉的狂欢中带离。如果我们的比喻被理解，我们就能轻易地将其中包含的规则应用于眼前的问题。那么，那个问题是什么呢？为什么当天父们竭尽全力保住家族继承人时，天主却常常让儿子和继承者在襁褓中被夺走？对那些如此发问的人，我们将用宴会的比喻来回答：即生命的宴席仿佛摆满了极其丰盛多样美味；而且必须注意到，正符合烹饪之道，并非所有菜肴都用享乐的蜂蜜调味，在某些情况下，生存被赋予了特别苦涩不幸的滋味：正如宴席艺术的大师们渴望用辛辣、咸或涩的菜肴来激发宾客的食欲。我说，生命并非在所有境遇中都如蜜般甘甜；其中有些境遇只有单纯的咸味，或者涩、酸、辛辣的味道渗透其中，使得浓汁难以入口；诱惑的杯中也盛满了各种饮料；有些因骄傲的错误产生浮夸虚荣的恶习；另一些引诱饮者做出鲁莽之举；而在其他情况下，它们引起呕吐，使多年非法所得之物羞耻地交出来。因此，为阻止一个过度狂欢的人继续留在这丰盛的筵席上，他早早地被从宾客中带走，从而得以逃脱那些无度放纵为饕餮之徒带来的灾祸。这就是我所说的完美眷顾的成就：不仅医治已犯的恶，而且防患于未然；而我们猜想，这就是新生婴儿死亡的原因。那按计划行事的一位，出于对个人的爱，移除了恶的原料，并对其预知已识别的性格，不给它时间通过实际邪恶的卓越表现来显示当邪恶倾向自由发挥时它是什么。同样，这生命筵席的安排者通过这样的处置，常常揭露贪爱钱财的\Quote{制约原因}的狡猾诡计，使这一恶习赤裸裸地暴露在阳光之下，毫无似是而非的托词，不再被任何误导性的屏障所掩盖。因为大多数人宣称他们放纵贪欲是为了使后代更加富有；但他们的恶习属于本性，并非由任何外在需要所引起，这在无子女者身上不可饶恕的贪婪中得到证明。许多没有继承人、也无希望有继承人的人，为他们辛苦积累的巨大财富，在心中养育了无数欲求而非子女，他们找不到渠道来输送这无可救药的疾病，尽管他们无法为这一缺点找到任何需要的借口。但有一些人，他们在尘世寄居期间，性情凶暴专横，成为各种情欲的奴隶，狂怒到疯狂，不放过任何最绝望的邪恶行为，是强盗、杀人犯、叛国者，更可憎的是，弑父、弑母、杀子，疯狂追求反常的性行为；假设这样的人在邪恶中老去；有人可能会问，这与我们先前探讨的结果如何一致？如果按照我们宴会的比喻，为避免其持续沉溺于生命的享乐，提早被带离，是出于眷顾而将其从狂欢中移除，那么对于某个具有那种性情的人，为何被允许继续他的狂欢直至老年，使他及其酒肉朋友都沉浸在他放荡的有害烟雾中？最后，你会问，为什么天主在其眷顾中，在一个人能在邪恶中臻于完善之前就将他从生命中带走，却让另一个人成长为如此怪物，以至于他从未出生倒好？对此，我们将向那些愿意善意接受的人提供如下理由：即那些生命良善之人的存在常常有益于其后代；在默感圣经中有数百处经文为此作证，清楚教导我们：天主对那些值得受其眷顾的人所施的温柔关怀，其后继者也共享；甚至对必定要过邪恶生活的人成为通往邪恶之路上的阻碍，也是一种好的结果。但鉴于我们在这一问题上理智不得不在黑暗中摸索，显然没有人能抱怨，如果其推测引领我们的思想得出各种结论。那么，不仅有人可以说，天主出于对某个家族奠基者的仁慈，将这家族中将要过邪恶生活的成员从邪恶生活中带走，而且，即使在某些早逝的情况下没有这样的先例，我们也不无理由推测，他们本会以比那些实际上已因邪恶而臭名昭著的人更为绝望的猛烈，陷入邪恶的生活。\par

我们从多处得知，没有天主，什么也不会发生；反之，凡认识到天主是理性、智慧、完美的美善和真理，且不能容忍那不善和与祂真理不符之事的人，都会承认天主的\OriginalTerm{眷顾安排}{dispensations}绝无任何偶然与混乱的成分。那么，婴儿的早逝究竟是归于上述原因，还是有某些超出这些的进一步原因，我们都应该承认，这些事的发生是为了最好的结果。我还有一个理由要提出，这是我从一位宗徒的智慧中学来的；也就是说，为什么有些以邪恶著称的人被容许在他们自选的路上继续活下去。在致罗马人的辩论中，宗徒较为详细地开展了这类想法，并针对由此必然引出的相反结论——即如果犯罪是天主的一种安排，那么罪人便不再能受正当责备，而且如果与掌管世界的天主意愿相悖，他根本就不会存在——他回应了这个结论，并借助一种更深的洞察解决了这个反驳。他告诉我们，天主在按每人的应得施予时，有时甚至容许邪恶最终产生善果。因此，祂容许例如\Person{法郎}出生并长成那样；目的是要那超越一切计算人数的大民族\Place{以色列}通过他的灾难得到教训。天主的\OriginalTerm{全能}{omnipotence}必须在每一方面都被承认；它有力量祝福配得者；它也不缺乏惩罚邪恶的能力；因此，既然需要将那特殊民族完全迁出\Place{埃及}，以免他们在误入歧途的生活方式中沾染埃及的罪恶，那反抗天主、声名狼藉的\Person{法郎}就在将要受益的民族的有生之年崛起并成熟，使以色列得以正确认识那\OriginalTerm{天主的双重能力}{two-fold energy of God}；他们在自己身上学到更仁慈的一面，通过看到天主对因邪恶而受鞭打者施以惩罚而学到更严厉的一面；因为天主以其至高的上智甚至能使邪恶与善合作。工匠（如果宗徒的论证可用我们的话来证实）——那工匠凭其技艺须将铁打造成日常用具，不仅需要那因其天然延展性而顺应其工艺的铁，而且无论铁有多么坚硬，无论它在火中多么难以软化，甚至因其金刚石般的阻力而无法锻造成任何有用器具，他的工艺也需要这铁的合作；他要用它作铁砧，在其上锤打柔软可锻的铁，使之成形为有用之物。但有人会说：\Quote{并非所有今生都像这样收获他们邪恶的果实，也并非所有生活有德的人都因他们的有德努力而获益；那么，我问，对于那些一直活到终了未受惩罚的人，他们存在的益处是什么呢？}我要提出一个超越所有人论证的理由来回应你的这个问题。伟大的\Person{达味}在某处宣称，义人的福乐有一部分将在于此：在目睹自己幸福的同时也看到被遗弃者的毁灭。他说：\BibleQuote{义人见复仇，必喜乐；他要以恶人的血洗手。}这并非因对受折磨者的痛苦而喜乐，而是那时最完全地认识到应得的德行奖赏的宏大。他用这话表明，将相反者与之对照，将是义人福乐的增添与加强。他说\BibleQuote{他以恶人的血洗手}，是想传达这样的思想：\Quote{他生活行为的清洁在恶人的毁灭中清楚地显明出来。}因为\Quote{洗}这个词语表达清洁的观念；但没有人是用血洗的，反而会被玷污；由此可见，是与更严厉的惩罚相比，才使德行的福乐显得更清晰。我们现在必须总结我们的论证，以便我们展开的思想更容易记忆。婴儿的早逝绝无任何暗示这样结束生命的人遭受某种严重不幸的成分，也不应等同于那些在此生通过各种德行净化自己的人的死亡；天主更具远见的\OriginalTerm{眷顾}{Providence}在那些将要过邪恶生活的人身上遏制了罪的巨大规模。一些恶人继续活着并不推翻我们所提出的这个理由；因为在他们身上，邪恶是出于对父母的仁慈而被阻止的；然而，对于那些父母从未赋予他们呼求天主能力的婴儿，这种伴随这种能力的天主仁慈形式并不会传递给他们的孩子；否则，那个因死亡而被阻止长大为恶的婴儿原会展现出比最臭名昭著的罪人更绝望的邪恶，因为他将不受阻碍。即使承认有些人已攀上罪恶的顶峰，宗徒的观点也为这个问题提供了安慰性的答案；因为那以智慧行事者知道如何借邪恶成就美善。更进一步，如果有些人罪恶昭彰，却从未像我们之前的比喻那样成为天主技艺用于善事的金属，那么，按先知的话所暗示，这构成了善人幸福的一种增加；这可以算作幸福中不可轻视的元素，也不能说是与天主的上智安排不相称的。\par

\SourceRecord{Translated by William Moore and Henry Austin Wilson.；Nicene and Post-Nicene Fathers, Second Series；Vol. 5.；Edited by Philip Schaff and Henry Wace.；Buffalo, NY: Christian Literature Publishing Co.,；1893.；<http://www.newadvent.org/fathers/2912.htm>.}
